\documentclass[../../../uem_phy_std_q.tex]{subfiles}

\begin{document}
    \begin{enumerate}
        \item 図のように，電池E（起電力$V_0$），電気抵抗R$_1$（抵抗値$2R$），%
            R$_2$（抵抗値$3R$）からなる回路がある．R$_2$の電位差（電圧）$V_2$を求めよ．\\
            \begin{minipage}{\linewidth}
                \centering
                    \includegraphics{\subfix{fig_uem_phy_std_em_cb_01/fig_uem_phy_std_em_cb_01_01_q.pdf}}
            \end{minipage}
            \vspace{1pt}
        \item 図のように，電池E（起電力$V_0$），電気抵抗R$_1$（抵抗値$2R$），R$_2$（抵抗値$R$），%
            R$_3$（抵抗値$2R$）からなる回路がある．%
            R$_2$を流れる電流の大きさ$i$を求めよ．\\
            \begin{minipage}{\linewidth}
                \centering
                    \includegraphics{\subfix{fig_uem_phy_std_em_cb_01/fig_uem_phy_std_em_cb_01_02_q.pdf}}
            \end{minipage}
            \vspace{1pt}
        \item 図のように，電池E$_1$（起電力$2V_0$），E$_2$（起電力$V_0$），%
            電気抵抗R$_1$（抵抗値$R$），R$_2$（抵抗値$R$），R$_3$（抵抗値$2R$）からなる回路がある．%
            AB間を流れる電流の大きさ$I$を求めよ．\\
            \begin{minipage}{\linewidth}
                \centering
                    \includegraphics{\subfix{fig_uem_phy_std_em_cb_01/fig_uem_phy_std_em_cb_01_03_q.pdf}}
            \end{minipage}
            \vspace{1pt}
        \item \needspace{10\baselineskip} 
            図のように，電池E（起電力$V_0$），電気抵抗R$_1$（抵抗値$2R$），R$_2$（抵抗値$R$），%
            R$_3$（抵抗値$R$），R$_4$（抵抗値$R$），R$_5$（抵抗値$R$）からなる回路がある．%
            AB間を流れる電流の大きさ$I$を求めよ．\\
            \begin{minipage}{\linewidth}
                \centering
                    \includegraphics{\subfix{fig_uem_phy_std_em_cb_01/fig_uem_phy_std_em_cb_01_04_q.pdf}}
            \end{minipage}            
    \end{enumerate}
\end{document}




